\chapter*{Where to find answers to key questions}
\addcontentsline{toc}{chapter}{Where to find answers to key questions}
\renewcommand{\currentnav}{Overview}
\setlength{\columnsep}{8mm}
\setlength{\columnseprule}{0.5pt}
\label{chap:Overview}

\begin{multicols}{2}

  \noindent This document serves two purposes. Firstly, it summarizes the quality of version~6.0\emph{x} Aura MLS \mbox{Level-2} and \mbox{Level-3} data. Secondly, it conveys important information to the community of MLS data users on how those data are to be read and interpreted.

  The MLS science team strongly encourages users of MLS data to thoroughly read the parts of the document relevant to their research interests. Chapter~\ref{chap:Essential} should be read by all, as it describes essential general information. Chapter~\ref{chap:Background} is considered background material that may be of interest to some data users. Chapter~\ref{chap:Level2} discusses individual MLS \mbox{Level-2} data (profiles of geophysical quantities along the orbit track) in detail, product by product.  Chapter~\ref{chap:Level3} describes the MLS \mbox{Level-3} data, which are averages of the quality-screened \mbox{Level-2} products in various coordinate systems (maps, zonal means, etc.).

  For convenience, this page provides answers, or pointers to answers, to anticipated questions.

  \subsection*{Where do I get \thisv MLS data?}
  All the MLS \mbox{Level-2} and \mbox{Level-3} data described here can be obtained from the NASA Goddard Space Flight Center Earth Science Data and Information Services Center (GES-DISC, \url{https://disc.gsfc.nasa.gov/}).  In addition, selected products are available on the Zenodo archive (\url{https://zenodo.org/communities/aura-mls}).

  \subsection*{What format are MLS data files in? }

  MLS \mbox{Level-2} data are in HDF-EOS version~5 format (see Section~\ref{sec:FilesAndFormats}, page~\pageref{sec:FilesAndFormats}). \mbox{Level-3} data are in netCDF4 (see Section~\ref{sec:L3files}, page~\pageref{sec:L3files}).

  \subsection*{Which MLS data points should be avoided?  How much should I trust
    the remainder?}

  These issues are described in Section~\ref{sec:StatusQuality} (starting on page~\pageref{sec:StatusQuality}), and on a product-by-product basis in Chapter~\ref{chap:Level2}. The key rules are:
  \begin{itemize}[leftmargin=2ex, itemsep=2pt]
    \item Only data within the appropriate pressure range (defined for each product in Chapter~\ref{chap:Level2}) should be used.
    \item Always consider the precision of the data, as reported in the \texttt{L2gpPrecision} field.
    \item Do not use any data points for which the precision is zero or a negative number, as that indicates poor information yield from MLS.
    \item Do not use data for any profile for which the corresponding entry in the integer \texttt{Status} array is an odd number.
    \item Data for profiles where \texttt{Status} is a non-zero even number should be approached with caution. See Section~\ref{sec:StatusQuality} on page~\pageref{sec:StatusQuality}, and the product descriptions in Chapter~\ref{chap:Level2} for details on how to interpret the \texttt{Status} information.
    \item Do not use any data for profiles for which the corresponding entry in the \texttt{Quality} array is \emph{smaller} than the threshold given in the section of Chapter~\ref{chap:Level2} describing your product of interest.
    \item Do not use any data for profiles for which the corresponding value in the \texttt{Convergence} array is \emph{larger} than the threshold given in the section of Chapter~\ref{chap:Level2} describing your product of interest.
    \item Some products require additional screening to remove biases or outliers; see specific sections of Chapter~\ref{chap:Level2} for details.
    \item Chapter~\ref{chap:Level2} also provides information on the accuracy of each product.
    \item Data users are strongly encouraged to contact the MLS science team (\href{mailto:mls-data@jpl.nasa.gov}{\texttt{mls-data@jpl.nasa.gov}}, or see the \hyperref[chap:Preface]{Preface} for alternate means of contact) to discuss their anticipated usage of the data and are always welcome to ask further data quality questions.
  \end{itemize}

  \subsection*{Why do some species abundances show negative values, and how do I
    interpret these?}

  Some of the MLS measurements have a poor signal-to-noise ratio for individual profiles. Radiance noise can naturally lead to some negative values for these species. It is critical to consider such values in scientific study. Any analysis that involves taking some form of average will exhibit a high bias if the points with negative mixing ratios are ignored.

\end{multicols}
